\section{Addendum: coordinated reheat delivery and delayed feedback}
\label{sec:oct27update}

\paragraph{Scope and historical record (27 October 2026).}
The preceding sections retain the measurements and interpretation of the
earlier manuscript. This addendum records subsequent experiments; its newer
results supersede statements about the current best controller or the absence
of an organizer evaluation, without replacing the earlier numerical record.
The previous source and compiled PDF are archived under
\texttt{reports/paper/history/}.

\subsection{Coordinating airflow and hot water}
The idle-handler airflow intervention, \texttt{d\_heat}, supplies minimum
fan pressure with a warm supply-air setpoint when a zone PI loop requests
heating. The handler itself cannot add heat. Consequently, an airflow override
with the hot-water pump stopped need not deliver the requested reheat.
The subsequent \texttt{e\_sync} controller also enables the hot-water loop
and restores its pump head when this override occurs, while retaining the
inherited demand-based water-temperature reset. It does not impose a new
constant hot-water temperature. The baseline for this comparison is
\texttt{h\_288}, the zone PI controller with the heating pressure clamp and
288-step cooling hold.

An initial five-scenario training screen reduced mean score from
$0.06122466$ for \texttt{h\_288} to $0.04162332$ for \texttt{d\_heat}
and $0.02252922$ for \texttt{e\_sync}. This screen included four high-comfort-loss
cases and a deep-cold control; it was not a full training-set result.
An initial three-scenario dev screen gave $0.04542558$ versus $0.01480846$.
The following completed sweeps include those cached full episodes and all
remaining scenarios, each with 8,640 simulation steps. Dev consists of
20 scenarios on four separately generated buildings and is a development
set, not an untouched final test set. The test split is not used here.

\begin{table}[htbp]
\centering
\small
\begin{tabular}{llrrrr}
\toprule
Split & Controller & Score & Comfort & Energy & Carbon \\
\midrule
Training ($n=30$) & \texttt{h\_288} & 0.02774466 & 0.01576395 & 0.00705920 & 0.00492151 \\
 & \texttt{e\_sync} & 0.01445394 & 0.00201534 & 0.00702192 & 0.00541668 \\
Dev ($n=20$) & \texttt{h\_288} & 0.03449141 & 0.02207958 & 0.00646678 & 0.00594505 \\
 & \texttt{e\_sync} & 0.01668440 & 0.00341916 & 0.00659122 & 0.00667402 \\
\bottomrule
\end{tabular}
\caption{Complete matched-cohort comparison. Lower score is better. The last
three columns are the weighted contributions to score, not raw physical units.
All 50 episodes completed without errors.}
\label{tab:syncfull}
\end{table}

The score reductions are $47.90\%$ on training and $51.63\%$ on dev.
On training, 24 scenarios improve, five regress and one is unchanged within
$10^{-10}$; on dev, 17 improve and three regress. The largest training
regression is approximately $0.00019424$, and the largest dev regression
$0.00012875$. Comfort improves substantially, while carbon cost rises on
both cohorts; this is not an improvement in every reward component.
The results and frozen run provenance are in
\texttt{experiments/runs/e\_sync\_coverage/}.

\paragraph{Organizer public score.}
The participant-provided Kaggle submission record shows scores
$0.06404$, $0.04932$, $0.03833$ and $0.01984$ for Versions 1--4 respectively.
Version 4 is the submitted \texttt{e\_sync} notebook. Its public score is
$48.24\%$ below Version 3. These are public leaderboard observations from
the submission record, not local scores or private-final-set measurements;
they are not pooled with the local cohorts in Table~\ref{tab:syncfull}.

\subsection{A subsequent setback experiment that did not advance}
A one-kelvin heating setback on top of \texttt{e\_sync},
\texttt{n\_sync\_s1}, was screened on two complete training episodes,
\texttt{fairbanks\_1} and \texttt{fairbanks\_7}. The corresponding mean score
increased from $0.03643457$ to $0.03694393$ ($+1.40\%$).
Unoccupied loss fell by $6.70\%$, but occupied loss rose by $11.30\%$.
The preselected gate required a lower mean training score, so this variant
did not advance to dev and did not replace Version 4. This negative result
is retained in \texttt{reports/NIGHT\_SETBACK\_REVIEW.md}.

\subsection{Observation semantics and a shadow predictor}
Source inspection and a three-step reproduction identify two constraints on
further feedback changes. First, the handler's
\texttt{supply\_air\_temperature\_sensor} reports its setpoint, rather than
the actual mixed/cooled supply temperature. In a controlled example with
return air at 295 K, outdoor air at 260 K, recirculation 0.7 and a 305 K
setpoint, the reported value is 305 K while actual supply is 284.5 K.
Second, the zone-temperature observation after a simulation step matches
the thermal state before that step, within $10^{-4}$ K in the reproduction.
The three-step diagnostic uses constant legal actions; it is not an evaluation
of controller performance. Internal thermal state is used only to diagnose
the delay, never as an input to a candidate controller.

With observation $o_t$ and action $a_t$, a causal predictor therefore estimates
the increment $o_{t+1}-o_t$ from observations available through $t$ and the
previous action $a_{t-1}$. A per-zone recursive ridge-regression shadow model
uses the zone temperature, its preceding observed increment, and the previous
zone, handler and water commands. Each target is scored before it is learned.
No future observations, future actions, building identifiers or hidden thermal
state enter its features.

On a 1,440-step \texttt{building\_train\_7/fairbanks\_7} trace of
\texttt{e\_sync}, with 288 prediction targets reserved for warmup, the scored
portion contains 1,151 steps across 36 zones. Holding the latest temperature
constant gives MAE 0.18872 K and RMSE 0.56355 K; constant-trend extrapolation
gives MAE 0.21891 K and RMSE 0.72990 K; the learned predictor gives
MAE 0.14335 K and RMSE 0.25009 K. This is evidence of predictive value on
one short trace, not evidence of lower closed-loop score or transfer to dev.
The trace's mean reward, $-0.01050882728$, reproduces the earlier diagnostic
rollout; it is not substituted for an 8,640-step episode score.

\subsection{Bounded closed-loop delay-compensation experiment}
The next experiment retains \texttt{e\_sync}'s aggregate handler and water
loops and applies the forecast only to zone PI feedback. Predictions are
bounded to $\pm1$ K after a 288-step warmup, and fixed blends of 0.5 and 1.0
are compared. A zero-blend control must reproduce \texttt{e\_sync}.
The protocol screens three full training episodes spanning deep-cold,
shoulder-season heating and hot-weather cooling. Only a candidate that passes
the training gate may advance to the predeclared dev screen and then broader
coverage. No notebook replaces Version 4 merely because prediction error falls.
The following automatically maintained record includes failures as well as
successful rollouts, and preserves earlier stages when later ones finish.

The training screen is fixed to \texttt{building\_train\_1/fairbanks\_1},
\texttt{building\_train\_7/fairbanks\_7} and
\texttt{building\_train\_8/phoenix\_4}. The zero-blend full-episode mean
reward must match the reference within $10^{-10}$. A nonzero blend must
lower the matched mean score without making any scenario more than $10\%$
worse. If both pass, only the lower training-mean candidate advances to
\texttt{building\_9001/fairbanks\_3},
\texttt{building\_9000/denver\_1},
\texttt{building\_9002/phoenix\_2} and
\texttt{building\_9003/singapore\_2}. The same mean and per-scenario
conditions gate the subsequent 30-training/20-dev sweep. Candidates and
thresholds are fixed before reading closed-loop outcomes.

\paragraph{Automatically recorded closed-loop results.}
Status: \texttt{complete}; stage: \texttt{finished}.
All entries below require complete 8,640-step episodes on matching scenarios.
\begin{center}\small
\begin{tabular}{llrrrr}\toprule
Cohort & Variant & $n$ & e\_sync & Candidate & Change \\\midrule
train screen & \texttt{r\_half} & 3 & 0.03207875 & 0.03274332 & +2.07\% \\
train screen & \texttt{r\_full} & 3 & 0.03207875 & 0.03198167 & -0.30\% \\
dev screen & \texttt{r\_full} & 4 & 0.02052920 & 0.02114155 & +2.98\% \\
\bottomrule\end{tabular}\end{center}
The training and dev screens are subsets of their respective full cohorts, not independent repetitions.
Completed new rollouts: 11. No candidate replaces Version 4 automatically.
The zero-blend full-episode control score is $0.01412923$; exact-reproduction gate: passed.
Experiment decision: dev gate failed; no full sweep or submission change.
The training-selected candidate is \texttt{r\_full}.


\subsection{Clarifications of the historical interpretation}
The numerical record above is retained, but several earlier statements need
qualification. The comfort price per kelvin includes each zone's occupancy
weight; it is not a universal unweighted per-zone charge. Tested setbacks
were worse than the flat-band reference, but their losses are not monotonic
in setback depth. The cooling hold has an empirical benefit, while a separate
fan-startup energy penalty or a cold-duct mechanism has not been established
by the inspected fan-power model. The heating delivery tradeoff is constrained
by the shipped 308.15 K VAV discharge limit and finite airflow; unchanged gas
use at arbitrarily low airflow cannot be assumed. Finally, a local demand or
climate estimate does not constitute a proven lower bound on the private
leaderboard score, and sharing a scalar parameter across buildings does not
by itself preclude overfitting. These qualifications change the interpretation
and strength of the claims, not the archived measurements.

\subsection{Follow-up: cooling-run persistence outside the PI window}
After delay compensation failed its dev gate, the next intervention changes
cooling-run persistence rather than temperature feedback. Version 4 can retain
a previously requested cooling run for 288 steps (24 hours), including hours
outside the PI and preconditioning window. The new candidates cap the remaining
timer outside that window to zero steps (\texttt{s\_night0}) or 12 steps
(one hour; \texttt{s\_night12}). Both preserve the full daytime 288-step hold,
the flat temperature bands, actual nighttime cooling/heating calls, and
the existing reheat airflow/water coordination. Thus this intervention does
not reintroduce the rejected temperature setback. It may reduce unoccupied
fan/chiller cost, but it may also lose useful overnight preconditioning;
the direction of the net score change requires full-episode measurement.
An earlier description incorrectly converted 288 steps to 12 hours; at the
configured 300 seconds per step the duration is 24 hours. This correction
changes no controller parameter or recorded episode result.

An unchanged 288-step cap is the reproduction control. The training and
dev screens, episode lengths and advancement gates are the same predeclared
sets and thresholds described above. A candidate is packaged as a separate
notebook only if complete coverage also improves both matched means and
does not worsen any individual case by more than $10\%$. Packaging additionally
checks standalone-versus-repository action parity on actual simulator
observations. No submission is sent to Kaggle automatically, and the Version 4
artifact and all previous experimental measurements are retained.

\paragraph{Automatically recorded closed-loop results.}
Status: \texttt{complete}; stage: \texttt{finished}.
All entries below require complete 8,640-step episodes on matching scenarios.
\begin{center}\small
\begin{tabular}{llrrrr}\toprule
Cohort & Variant & $n$ & e\_sync & Candidate & Change \\\midrule
train screen & \texttt{s\_night0} & 3 & 0.03207875 & 0.03156326 & -1.61\% \\
train screen & \texttt{s\_night12} & 3 & 0.03207875 & 0.03192036 & -0.49\% \\
dev screen & \texttt{s\_night0} & 4 & 0.02052920 & 0.02103063 & +2.44\% \\
\bottomrule\end{tabular}\end{center}
The training and dev screens are subsets of their respective full cohorts, not independent repetitions.
Completed new rollouts: 11. No candidate replaces Version 4 automatically.
The unchanged-hold full-episode control score is $0.01412923$; exact-reproduction gate: passed.
Experiment decision: dev gate failed; no full sweep or submission change.
The training-selected candidate is \texttt{s\_night0}.


\subsection{Exploratory follow-up of the one-hour cooling hold}
The training-selected zero-hold variant failed the four-case dev screen:
its matched mean increased from $0.02052920$ to $0.02103063$ ($+2.44\%$),
with a worst individual regression of $17.59\%$. The original experiment
therefore stopped as specified and produced no replacement submission.
The alternative one-hour variant, \texttt{s\_night12}, had also passed the
training gate ($0.03207875$ to $0.03192036$, or $-0.49\%$) but was not
evaluated on dev under that original single-candidate protocol.

The follow-up now evaluates this remaining candidate without changing its
parameters. The three training episodes and unchanged-control episode are
explicitly reused after checking source hashes, variant parameters and
original result fingerprints; they are not additional independent evidence.
The same four dev cases remain development data, and their reuse after the
previous candidate's failure means this is exploratory model selection,
not an untouched confirmation set. The separate test split remains unused.
Only a passing dev screen advances to full matched training/dev coverage and
standalone notebook verification, under the same mean and worst-case gates.
Earlier candidate outcomes and all historical measurements remain retained.

\paragraph{Automatically recorded closed-loop results.}
Status: \texttt{complete}; stage: \texttt{finished}.
All entries below require complete 8,640-step episodes on matching scenarios.
This is an exploratory follow-up after the zero-hold candidate failed dev; training results and the control are reused, not new independent episodes.
\begin{center}\small
\begin{tabular}{llrrrr}\toprule
Cohort & Variant & $n$ & e\_sync & Candidate & Change \\\midrule
train screen & \texttt{s\_night12} & 3 & 0.03207875 & 0.03192036 & -0.49\% \\
dev screen & \texttt{s\_night12} & 4 & 0.02052920 & 0.02119149 & +3.23\% \\
\bottomrule\end{tabular}\end{center}
The training and dev screens are subsets of their respective full cohorts, not independent repetitions.
Completed new rollouts: 8. No candidate replaces Version 4 automatically.
The unchanged-hold full-episode control score is $0.01412923$; exact-reproduction gate: passed.
Experiment decision: dev gate failed; no full sweep or submission change.
The training-selected candidate is \texttt{s\_night12}.


\subsection{Pump-head reset with a comfort and valve-demand guard}
The one-hour hold follow-up also failed its dev gate: the mean increased to
$0.02119149$ from the matched $0.02052920$ ($+3.23\%$), with a worst individual
regression of $18.89\%$. Thus neither shortened-hold variant replaces Version 4.
The next experiment retains Version 4's original temperature feedback, PI gains,
288-step cooling hold and airflow/water interlock. It instead modifies the
final hot-water pump-pressure request, after all coordination overrides.

Reducing differential pressure lowers both pump electricity and water flow,
so it can reduce heating delivery. The earlier unconditional pump-cap failure
on the older fan controller is retained as a negative result, not dismissed.
The new experiment includes an unconditional cap on \texttt{e\_sync}
(\texttt{u\_pump\_cap}) as an ablation and a guarded version
(\texttt{u\_pump\_guard}). Both cap the normalized pressure action at zero,
which corresponds to 10 bar under the shipped 0--20 bar normalization.
The guarded version preserves the reference request whenever any commanded
open reheat valve serves a zone at or below the comfort lower bound plus
0.5 K, has opening at least 0.8, or has a non-finite temperature observation.
No cap is applied when the pump is off or all reheat valves are closed;
zero-flow pump electricity was already zero in the model.

This is a proposed tradeoff, not a guarantee of unchanged heat delivery.
In particular, the guard uses delayed observations and cannot prove that
the current hidden thermal state is safe. The 0.5 K reserve, 0.8 valve threshold
and normalized cap are fixed before this experiment's outcomes. Only allowed
observations and final requested actions enter the guard; no true thermal state,
building identifier or weather label is used for control.

An uncapped control must reproduce one full reference episode. Both candidates
are then screened on the same three training cases. Every training-passing
candidate is evaluated on the same four dev cases; among those passing dev,
the lower dev mean selects the full-coverage candidate. At each gate the matched
mean must fall and no individual score may rise by more than $10\%$.
The dev set is reused development data for exploratory selection, not an
untouched confirmation set. The test split remains unused. Packaging requires
complete 30-training/20-dev coverage passing both gates and standalone action
parity verification. Failed or short episodes cannot qualify.

\paragraph{Automatically recorded closed-loop results.}
Status: \texttt{complete}; stage: \texttt{finished}.
All entries below require complete 8,640-step episodes on matching scenarios.
\begin{center}\small
\begin{tabular}{llrrrr}\toprule
Cohort & Variant & $n$ & e\_sync & Candidate & Change \\\midrule
train screen & \texttt{u\_pump\_cap} & 3 & 0.03207875 & 0.04016067 & +25.19\% \\
train screen & \texttt{u\_pump\_guard} & 3 & 0.03207875 & 0.03194406 & -0.42\% \\
dev screen & \texttt{u\_pump\_guard} & 4 & 0.02052920 & 0.02060415 & +0.37\% \\
\bottomrule\end{tabular}\end{center}
The training and dev screens are subsets of their respective full cohorts, not independent repetitions.
Completed new rollouts: 11. No candidate replaces Version 4 automatically.
The unchanged-pump full-episode control score is $0.01412923$; exact-reproduction gate: passed.
Experiment decision: all candidate dev gates failed; no submission change.


\subsection{Follow-up: restrict pump economy to heating commands}
\paragraph{Update dated 28 October 2026.}
The preceding pump experiment completed 11 valid episodes. An unconditional
cap increased the three-case training mean by $25.19\%$. The guarded cap
reduced that training mean by $0.42\%$ but increased the four-case dev mean
from $0.02052920$ to $0.02060415$ ($+0.37\%$); neither produced a replacement
notebook. The largest dev regression was $3.11\%$ in Singapore, primarily
from the comfort term. This motivates a narrower intervention rather than
assuming the pump reduction is harmless during cooling.

Both new arms retain the same 0.5 K reserve, 0.8 valve-demand threshold and
normalized-zero cap, but refuse pressure economy whenever any running
handler's temperature command is nonpositive. This is a test on the requested
control direction, not a reading of the misleading supply-temperature sensor
and not a claim that the handler itself can heat. The broader arm,
\texttt{v\_pump\_warm}, may cap during any eligible all-heating-command step.
The stricter arm, \texttt{v\_pump\_heat}, additionally requires that
\texttt{e\_sync} added a reheat-airflow override on this same step. The override
counter counts command overrides, not physical fan-start transitions.
These gates use the controller's own state and allowed observations, without
building identifiers, climate labels or future temperatures.

The new experiment repeats the unchanged-control episode and both three-case
training screens. All training-passing arms receive the same four development
cases; the lower passing dev mean selects full 30-training/20-dev coverage.
The previously stated mean and worst-case gates, full episode lengths and
standalone notebook verification still apply. Repeated use of dev is explicitly
exploratory model development; it is not an untouched confirmation set.
The test split remains unused. The scheduler waits for available RAM and spaces
interpreter launches to avoid simultaneous initialization peaks.

\paragraph{Automatically recorded closed-loop results.}
Status: \texttt{complete}; stage: \texttt{finished}.
All entries below require complete 8,640-step episodes on matching scenarios.
\begin{center}\small
\begin{tabular}{llrrrr}\toprule
Cohort & Variant & $n$ & e\_sync & Candidate & Change \\\midrule
train screen & \texttt{v\_pump\_warm} & 3 & 0.03207875 & 0.03203897 & -0.12\% \\
train screen & \texttt{v\_pump\_heat} & 3 & 0.03207875 & 0.03203897 & -0.12\% \\
dev screen & \texttt{v\_pump\_warm} & 4 & 0.02052920 & 0.02062192 & +0.45\% \\
dev screen & \texttt{v\_pump\_heat} & 4 & 0.02052920 & 0.02062192 & +0.45\% \\
\bottomrule\end{tabular}\end{center}
The training and dev screens are subsets of their respective full cohorts, not independent repetitions.
Completed new rollouts: 15. No candidate replaces Version 4 automatically.
The unchanged-pump full-episode control score is $0.01412923$; exact-reproduction gate: passed.
Experiment decision: all candidate dev gates failed; no submission change.


\subsection{Residual-comfort diagnosis before an earlier-prestart sweep}
Both heating-command pump variants produced identical scores on the seven
screening cases: the training mean fell by $0.12\%$, but the four-case dev mean
rose by $0.45\%$. The 15-episode experiment therefore stopped without a new
submission. The next proposed hypothesis was to extend the existing one-hour
PI preconditioning period to two or three hours, without changing other V4
parameters. Before running those candidates, a full reference diagnostic was
required to show that the first two scheduled comfort hours accounted for
at least $25\%$ of comfort loss. This heuristic concentration gate was fixed
before inspecting that diagnostic's result.

On \texttt{building\_train\_7/phoenix\_3}, V4's largest residual training
comfort contributor, the 8,640-step diagnostic reproduced score
$0.01620891807$. Its hourly grouping needed a clock correction: converting
protobuf timestamps to US/Pacific shifted bins relative to the naive simulation
wall clock used by observations and the comfort schedule. The initial diagnostic
is retained as an audit artifact; the corrected repeat uses observed policy
clock features and reproduces the same full score. Neither simulator nor
controller timing was modified.

The corrected first-two-hour comfort share is $12.00\%$, below both the
predeclared $25\%$ gate and the $15.38\%$ share of a uniform 13-hour window.
The proposed earlier-prestart sweep was therefore not run; no performance
claim is made for the registered two- and three-hour candidates. Instead,
$96.30\%$ of this diagnostic's comfort loss comes from zones above their
cooling bounds. Approximately $92.47\%$ of total comfort loss co-occurs with
a lagged hot-zone observation, positive cooling PI demand, a running cooling
handler and a pressure command at or below the shipped 200 Pa floor.
Approximately $58.07\%$ co-occurs with such a zone's PI demand at least 0.8.
These are correlations within one training episode, not causal effects or
global-cohort estimates. Directional allocation uses the reward function's own
zone-productivity curve, occupancy and reward information, strictly offline.
No true reward-state temperatures are supplied to a submitted controller.

\paragraph{Automatically recorded closed-loop results.}
Status: \texttt{complete}; stage: \texttt{hypothesis\_not\_supported}.
All entries below require complete 8,640-step episodes on matching scenarios.
\begin{center}\small
\begin{tabular}{llrrrr}\toprule
Cohort & Variant & $n$ & e\_sync & Candidate & Change \\\midrule
\bottomrule\end{tabular}\end{center}
The training and dev screens are subsets of their respective full cohorts, not independent repetitions.
Completed recorded rollouts: 1. No candidate replaces Version 4 automatically.
The hourly diagnostic is reused as the control, not an additional independent repetition.
First two scheduled hours account for 12.00\% of the diagnostic comfort loss.
Experiment decision: first-two-hour comfort share below 25%; do not run the proposed prestart sweep.


\subsection{Bounded pressure response to saturated hot-zone PI demand}
The next experiment uses that residual-comfort diagnosis to test cooling
airflow authority directly. It retains V4's PI gains, band centre, one-hour
prestart, pump requests, hot-water reset and full 288-step cooling hold.
Only the pressure of an already running cooling handler may be increased.
At least one of its zones must be above its cooling bound with positive
zone PI demand at least 0.8, during the PI window.

Let $d$ be the largest positive upper-band excursion among those eligible
zones and $f=\min(1,d/(1\,\mathrm{K}))$. Under the shipped normalized pressure
range, the candidate request is
$p^*=-0.8+f(p_{\max}+0.8)$, applied only if larger than the reference command.
The candidates use $p_{\max}=-0.6$ or $-0.4$, corresponding to 400 or 600 Pa
under the shipped 0--2,000 Pa normalization; $-0.8$ represents the model's
200 Pa running floor. A missing or non-finite outdoor temperature vetoes
the modification. A mixed-air proxy using the requested recirculation capped
at 0.7 and the mean observed zone return temperature also vetoes it when below
the comfort lower bound. This proxy is approximate and uses delayed zone
sensors; it is not a direct actual-supply-temperature measurement.

More pressure increases fan power and may overcool other zones, so the
one-case diagnostic does not establish a benefit. The unchanged control and
both candidates are evaluated on the three largest residual-training-comfort
cases (\texttt{phoenix\_3}, \texttt{atlanta\_5}, \texttt{denver\_5}) and a
deep-cold \texttt{fairbanks\_1} control. Training-passing arms receive the
three high-comfort-loss dev cases on \texttt{building\_9002}
(Phoenix, Atlanta, Singapore) and a Fairbanks heating control on
\texttt{building\_9001}. These deliberately targeted, repeatedly used
development cohorts are not representative global estimates or independent
confirmation sets. The same full-episode, matched-mean, worst-case and
complete-coverage packaging gates remain in force. The test split is unused.

\paragraph{Automatically recorded closed-loop results.}
Status: \texttt{complete}; stage: \texttt{finished}.
All entries below require complete 8,640-step episodes on matching scenarios.
\begin{center}\small
\begin{tabular}{llrrrr}\toprule
Cohort & Variant & $n$ & e\_sync & Candidate & Change \\\midrule
train screen & \texttt{cpi\_400} & 4 & 0.02564186 & 0.02541995 & -0.87\% \\
train screen & \texttt{cpi\_600} & 4 & 0.02564186 & 0.02544590 & -0.76\% \\
dev screen & \texttt{cpi\_400} & 4 & 0.01684390 & 0.01741959 & +3.42\% \\
dev screen & \texttt{cpi\_600} & 4 & 0.01684390 & 0.01766523 & +4.88\% \\
\bottomrule\end{tabular}\end{center}
The training and dev screens are subsets of their respective full cohorts, not independent repetitions.
Completed new rollouts: 17. No candidate replaces Version 4 automatically.
The unchanged-pump full-episode control score is $0.01620892$; exact-reproduction gate: passed.
Experiment decision: all candidate dev gates failed; no submission change.


\subsection{Flow compensation when raising a shared AHU pressure}
The bounded-pressure experiment completed 17 valid episodes. The 400 Pa arm
reduced the four-case training mean by $0.87\%$ but increased the four-case dev
mean by $3.42\%$; the 600 Pa arm gave $-0.76\%$ and $+4.88\%$, respectively.
Both were rejected before full coverage or packaging. On dev Phoenix,
comfort loss fell, whereas on Atlanta and Singapore it rose. This motivates
testing cross-zone airflow coupling, without claiming that this pattern
already proves the causal mechanism.

In the shipped VAV model, nominal airflow satisfies
$Q_i\propto d_i\sqrt{P}$, where $d_i$ is damper opening and $P$ is the
handler's effective running pressure. Raising a shared pressure therefore
increases airflow in non-target rooms as well as in the hot target.
The new arms retain the previous 400 Pa response and its eligibility/mixing
gates. For a non-target zone, they request
\[
 d_i'=\max\!\left(d_{\min},\ d_i\sqrt{P_{\mathrm{old}}/P_{\mathrm{new}}}\right),
 \qquad d_{\min}=0.1,
\]
using effective reference opening/pressure after the shipped ventilation and
200 Pa floors. An eligible saturated hot target retains its original opening.
Reheat-valve commands, hot-water commands, SAT and recirculation requests are
not rewritten. Above the damper floor this preserves nominal non-target flow
in that model; at the floor exact compensation is impossible. It also does
not guarantee identical heat transfer, since return/mixed temperatures and
thermal feedback remain coupled.

\texttt{cfb\_comp} applies the clipped compensation. \texttt{cfb\_guard}
additionally vetoes the entire handler modification if a non-target zone at
or below the lower comfort bound plus 0.5 K would need an opening below the
0.1 floor. A missing non-target damper command also vetoes the modification.
Updates are collected before being applied, so a veto cannot leave partially
rewritten dampers. The guard uses delayed observations and is not a safety
guarantee. Its parameters are fixed before this experiment's outcomes.

The unchanged control, four training cases, four repeatedly used dev cases,
selection gates and full-coverage packaging gates remain those of the preceding
pressure experiment. The uncompensated 400 Pa results remain an archived
ablation comparator on matching cases, not additional independent evidence.
The original V4 controller and its submitted artifact are retained; no new
public-score improvement is claimed from these targeted development screens.
The separate test split remains unused.

\paragraph{Automatically recorded closed-loop results.}
Status: \texttt{complete}; stage: \texttt{finished}.
All entries below require complete 8,640-step episodes on matching scenarios.
\begin{center}\small
\begin{tabular}{llrrrr}\toprule
Cohort & Variant & $n$ & e\_sync & Candidate & Change \\\midrule
train screen & \texttt{cfb\_comp} & 4 & 0.02564186 & 0.02525925 & -1.49\% \\
train screen & \texttt{cfb\_guard} & 4 & 0.02564186 & 0.02525925 & -1.49\% \\
dev screen & \texttt{cfb\_comp} & 4 & 0.01684390 & 0.01720231 & +2.13\% \\
dev screen & \texttt{cfb\_guard} & 4 & 0.01684390 & 0.01713035 & +1.70\% \\
\bottomrule\end{tabular}\end{center}
The training and dev screens are subsets of their respective full cohorts, not independent repetitions.
Completed new rollouts: 17. No candidate replaces Version 4 automatically.
The unchanged-pump full-episode control score is $0.01620892$; exact-reproduction gate: passed.
Experiment decision: all candidate dev gates failed; no submission change.


\subsection{Full-development calibration of zone PI gains}
The flow-compensated arms completed 17 valid episodes. Both reduced the
four-case training mean by $1.49\%$, but the unguarded and guarded arms
increased the four-case dev mean by $2.13\%$ and $1.70\%$, respectively.
They were rejected and no replacement notebook was produced. These results
remain in the preceding table.

The next study returns to the unchanged \texttt{e\_sync} plant and calibrates
only the per-zone PI coefficients. Earlier soft/stiff PI probes changed both
gains on an older controller while retaining a fixed integral-state limit.
For example, halving $K_i$ at a fixed limit of 200 also halves the maximum
integral contribution from 1 to 0.5. That confounds response speed with
available integral authority. Here every candidate sets
$I_{\max}=1/K_i$, preserving $K_i I_{\max}=1$, while all handler, pump, water,
band, one-hour prestart and 288-step cooling-hold settings retain V4 values.

\begin{center}\small
\begin{tabular}{lrrr}\toprule
Candidate & $K_p$ & $K_i$ & $I_{\max}$ \\\midrule
V4 reproduction & 0.10 & 0.0050 & 200 \\
Slower integral & 0.10 & 0.0025 & 400 \\
Faster integral & 0.10 & 0.0100 & 100 \\
Softer proportional & 0.05 & 0.0050 & 200 \\
Firmer proportional & 0.20 & 0.0050 & 200 \\
\bottomrule\end{tabular}\end{center}

A fresh complete V4-control episode must reproduce the cached reference
within $10^{-10}$ mean reward. The four one-coordinate neighbors are then
evaluated on \emph{all 20 dev scenarios}, rather than a small targeted screen.
Candidates with lower full-dev means are ranked and checked on all 30 training
scenarios in that order. If an earlier-ranked candidate fails the training
mean check, its complete results remain recorded and the next candidate is
checked automatically. If no first-batch candidate improves both full means,
the fixed second batch tests the four joint corners
$K_p\in\{0.05,0.20\}$ and $K_i\in\{0.0025,0.0100\}$, again with
$I_{\max}=1/K_i$. The first candidate passing both complete means is packaged
and standalone-action verified. This bounded search does not claim a global
optimum or require the user to restart after each failed neighbor.

For this new prospective study, selection uses the competition's mean-score
objective, with a $10^{-10}$ absolute numerical tolerance. The previous
heuristic $10\%$ per-scenario cutoff is not imposed; regression counts and the
largest individual increases are still reported. Previous failed studies are
not reclassified: their dev means were worse. Dev is repeatedly used tuning
data and the training set is a published consistency check, so neither is an
untouched confirmation set. The separate test split remains unused, and no
local advantage guarantees a lower organizer public/private score.
Source/config/reference hashes, installed dependency versions, a duplicate-run
lock, available-memory checks and per-rollout checkpoints accompany the study.
All earlier measurements and the V4 artifact remain retained.

\paragraph{Automatically recorded closed-loop results.}
Status: \texttt{complete}; stage: \texttt{finished}.
All entries below require complete 8,640-step episodes on matching scenarios.
\begin{center}\small
\begin{tabular}{llrrrr}\toprule
Cohort & Variant & $n$ & e\_sync & Candidate & Change \\\midrule
full dev & \texttt{pi\_cal\_slow\_i} & 20 & 0.01668440 & 0.01793284 & +7.48\% \\
full dev & \texttt{pi\_cal\_fast\_i} & 20 & 0.01668440 & 0.01656292 & -0.73\% \\
full dev & \texttt{pi\_cal\_soft\_p} & 20 & 0.01668440 & 0.01707339 & +2.33\% \\
full dev & \texttt{pi\_cal\_firm\_p} & 20 & 0.01668440 & 0.01636229 & -1.93\% \\
full dev & \texttt{pi\_cal\_soft\_slow} & 20 & 0.01668440 & 0.01939039 & +16.22\% \\
full dev & \texttt{pi\_cal\_soft\_fast} & 20 & 0.01668440 & 0.01669641 & +0.07\% \\
full dev & \texttt{pi\_cal\_firm\_slow} & 20 & 0.01668440 & 0.01661863 & -0.39\% \\
full dev & \texttt{pi\_cal\_firm\_fast} & 20 & 0.01668440 & 0.01670230 & +0.11\% \\
full training & \texttt{pi\_cal\_fast\_i} & 30 & 0.01445394 & 0.01502695 & +3.96\% \\
full training & \texttt{pi\_cal\_firm\_p} & 30 & 0.01445394 & 0.01469519 & +1.67\% \\
full training & \texttt{pi\_cal\_firm\_slow} & 30 & 0.01445394 & 0.01459717 & +0.99\% \\
\bottomrule\end{tabular}\end{center}
These are complete matched cohorts. Dev is repeatedly used model-selection data, not untouched validation. Individual regressions are reported without a hard cutoff.
Completed new rollouts: 251. No candidate replaces Version 4 automatically.
The unchanged-PI full-episode control score is $0.01620892$; exact-reproduction gate: passed.
Experiment decision: no tested PI neighbor improved both full means; keep V4.


\subsection{Geometry-scaled zone PI: prospective experiment, 29 October}
The scalar PI study completed 251 valid episodes. Three candidates reduced
the full dev mean but increased the full training mean: firmer proportional
gain by $1.669\%$, faster integral gain by $3.964\%$, and firmer proportional
with slower integral gain by $0.991\%$. None replaced V4; their preceding
tables and the public V4 score remain unchanged.

The next hypothesis accounts for unequal zone sizes. Current VAV maximum
airflow is uniform, whereas thermal mass depends on geometry. For room area
$A$ measured in control volumes, define
$s=\operatorname{clip}((A/2484)^p,0.5,2)$.
The reference 2484 is the median across 323 zones in the six training
buildings only. Three fixed rules test $(p,K_p,K_i,I_{\max})$ equal to
$(0.5,0.1s,0.005,200)$, $(0.5,0.1s,0.005s,200/s)$, and
$(1,0.1s,0.005,200)$. All retain maximum integral contribution one,
conditional integration and every other V4 plant/schedule setting.

The submitted policy derives areas from participant-facing \texttt{zone\_map}
metadata with a NumPy/Python four-connected component scan. An offline check
confirms component areas and SDK room/VAV numbering on all ten current
training/dev buildings; dev areas do not set the training reference. Runtime
uses no simulator state, reward or building/weather identity. Missing maps or
inconsistent VAV IDs fall back to V4. A same-count label permutation on another
SDK version would not be detected by this fallback; compatibility beyond
the verified SDK/maps is not established.

After a full disabled-geometry control reproduces V4, all three candidates
receive all 20 dev episodes. Dev-improving rules receive all 30 training
episodes in dev-rank order until both full means improve, followed by
standalone parity verification and notebook packaging. The bounded study
uses at most 151 full episodes before standalone checks. It uses the same
mean-only objective and numerical tolerance as the preceding study, reports
individual regressions, and leaves the test split unused. Repeated dev tuning
and a local improvement do not establish a better organizer score.

\paragraph{Automatically recorded closed-loop results.}
Status: \texttt{complete}; stage: \texttt{finished}.
All entries below require complete 8,640-step episodes on matching scenarios.
\begin{center}\small
\begin{tabular}{llrrrr}\toprule
Cohort & Variant & $n$ & e\_sync & Candidate & Change \\\midrule
full dev & \texttt{geo\_p\_sqrt} & 20 & 0.01668440 & 0.01621940 & -2.79\% \\
full dev & \texttt{geo\_pi\_sqrt} & 20 & 0.01668440 & 0.01609201 & -3.55\% \\
full dev & \texttt{geo\_p\_linear} & 20 & 0.01668440 & 0.01614269 & -3.25\% \\
full training & \texttt{geo\_pi\_sqrt} & 30 & 0.01445394 & 0.01415713 & -2.05\% \\
\bottomrule\end{tabular}\end{center}
These are complete matched cohorts. Dev is repeatedly used model-selection data, not untouched validation. Individual regressions are reported without a hard cutoff.
Completed new rollouts: 91. No candidate replaces Version 4 automatically.
The unchanged-PI full-episode control score is $0.01620892$; exact-reproduction gate: passed.
The selected candidate is \texttt{geo\_pi\_sqrt}.


\subsection{V6: refining the geometry exponent after V5 selection}
The geometry study completed 91 full episodes without rollout errors. Its
selected \texttt{geo\_pi\_sqrt} policy (submitted as V5, organizer score pending)
has a full training mean of $0.0141571250$ against V4's $0.0144539435$
($-2.05\%$), and full dev mean $0.0160920115$ against $0.0166844015$
($-3.55\%$). It improves 27 of 30 training and 17 of 20 dev scenarios;
the largest relative regression is $3.12\%$ on training and $4.80\%$ on dev.
Standalone verification matched 1,152 action pairs across two scenarios.
All preceding measurements remain retained.

V5's dev weighted comfort penalty falls from $0.0034191613$ to
$0.0027456800$ (approximately $19.70\%$), while weighted energy and carbon
rise slightly. The next bounded study refines only the geometry exponent:
$s=\operatorname{clip}((A/2484)^p,0.5,2)$ with
$p\in\{0.25,0.75,1.0\}$ around V5's $p=0.5$.
Every candidate uses $K_p=0.1s$, $K_i=0.005s$ and $I_{\max}=200/s$.
The geometry implementation, training-only reference and all plant/schedule
settings are unchanged. Weaker and stronger sensitivities are hypotheses;
the successful first exponent is not assumed optimal.

A fresh full V5 control must reproduce its reference. All three candidates
receive all 20 dev scenarios; those improving both V5 and V4 full-dev means
receive all 30 training scenarios in dev-rank order, continuing after a failed
training check. Packaging and standalone verification require both complete
cohort means to improve on \emph{both} references by more than $10^{-10}$.
The maximum is 151 new full episodes before standalone verification.
Regression details remain visible; repeated dev tuning is exploratory and the
test split remains unused. V4/V5 artifacts and earlier tables remain intact.
Only organizer evaluation can establish an improvement over their Kaggle scores.

\paragraph{Automatically recorded closed-loop results.}
Status: \texttt{complete}; stage: \texttt{finished}.
All entries below require complete 8,640-step episodes on matching scenarios.
\begin{center}\small
\begin{tabular}{llrrrr}\toprule
Cohort & Variant & $n$ & V5 & Candidate & Change \\\midrule
full dev & \texttt{geo6\_quarter} & 20 & 0.01609201 & 0.01623829 & +0.91\% \\
full dev & \texttt{geo6\_threequarter} & 20 & 0.01609201 & 0.01604957 & -0.26\% \\
full dev & \texttt{geo6\_linear} & 20 & 0.01609201 & 0.01616189 & +0.43\% \\
full training & \texttt{geo6\_threequarter} & 30 & 0.01415713 & 0.01407835 & -0.56\% \\
\bottomrule\end{tabular}\end{center}
These are complete matched cohorts. Dev is repeatedly used model-selection data, not untouched validation. Individual regressions are reported without a hard cutoff.
Completed new rollouts: 91. V5 remains the submitted reference; organizer improvement is not established.
The unchanged-V5 full-episode control score is $0.01564186$; exact-reproduction gate: passed.
The selected candidate is \texttt{geo6\_threequarter}.


\subsection{Organizer feedback for V5, recorded 30 October}
The user's Kaggle result screenshot reports V5 at $0.01996$, compared with
V4 at $0.01984$. Lower is better: V5 is worse by $0.00012$, approximately
$0.60\%$. V4 therefore remains the best organizer-scored submission among
these two versions. The earlier pending-score statements describe the state
when their experiments were recorded; this paragraph adds the later result
without replacing the local measurements or earlier tables.

V5's lower complete training/dev means did not translate into an organizer
improvement. This is direct evidence that repeatedly tuned local cohorts
cannot establish the organizer ranking. The prospective V6 experiment retains
its frozen protocol and local comparisons with both references; even a local
pass will remain an experimental submission requiring organizer evaluation.

\subsection{V6 completion and prospective V7 split geometry gains}
The V6 study completed 91 full episodes with no rollout errors. Selected
\texttt{geo6\_threequarter} reached training mean $0.0140783514$ and dev mean
$0.0160495740$, reductions of $0.556\%$ and $0.264\%$ relative to V5.
It improved 23 of 30 training and 8 of 20 dev cases; the largest relative
regressions versus V5 were $2.94\%$ and $3.65\%$, respectively.
Standalone verification matched 1,152 action pairs across two scenarios.
V6 was submitted; its organizer score is pending at this experiment's start.
V4 remains the best known organizer-scored version. Earlier tables remain intact.

The next hypothesis separates the room-area sensitivity of instantaneous
proportional response from accumulated integral response. V6 reduced the dev
weighted comfort penalty from V5's $0.0027456800$ to $0.0026861601$, while
weighted carbon rose from $0.0067092221$ to $0.0067357395$. These measurements
motivate a test, not a physical identification claim. Define
$s_p=\operatorname{clip}((A/2484)^p,0.5,2)$ and
$s_i=\operatorname{clip}((A/2484)^i,0.5,2)$, with gains
$K_p=0.1s_p$, $K_i=0.005s_i$ and raw integral limit $200/s_i$.
The three fixed candidates use $(p,i)=(0.75,0.5),(0.75,1.0),(0.5,0.75)$;
the reproduction control uses $(0.75,0.75)$, exactly V6's rule.
The underlying PI conditional integration and actuator law are reused.
All geometry mapping assumptions and V6 plant/schedule settings remain unchanged.

After a complete reproduction control, all candidates receive all 20 dev
scenarios. Candidates reducing the mean against V6, V5 \emph{and} V4 receive
all 30 training scenarios in dev-rank order; packaging and standalone checks
require both full means to beat all three references by more than $10^{-10}$.
The bounded maximum is 151 new episodes before standalone checks.
Individual regressions are reported without a hard cutoff. Repeated local
tuning is exploratory; the test split stays unused and organizer superiority
requires a Kaggle result. Source/config/reference hashes, available-memory
and disk checks, checkpoints and a duplicate-run lock accompany the study.
New PDF builds retain the previous readable PDF unless compilation succeeds.

\paragraph{Automatically recorded closed-loop results.}
Status: \texttt{complete}; stage: \texttt{finished}.
All entries below require complete 8,640-step episodes on matching scenarios.
\begin{center}\small
\begin{tabular}{llrrrr}\toprule
Cohort & Variant & $n$ & V6 & Candidate & Change \\\midrule
full dev & \texttt{geo7\_soft\_i} & 20 & 0.01604957 & 0.01605772 & +0.05\% \\
full dev & \texttt{geo7\_strong\_i} & 20 & 0.01604957 & 0.01603589 & -0.09\% \\
full dev & \texttt{geo7\_soft\_p} & 20 & 0.01604957 & 0.01613853 & +0.55\% \\
full training & \texttt{geo7\_strong\_i} & 30 & 0.01407835 & 0.01422959 & +1.07\% \\
\bottomrule\end{tabular}\end{center}
These are complete matched cohorts. Dev is repeatedly used model-selection data, not untouched validation. Individual regressions are reported without a hard cutoff.
Completed new rollouts: 91. V6 remains the latest submitted candidate; organizer improvement is not established.
The unchanged-V6 full-episode control score is $0.01543014$; exact-reproduction gate: passed.
Experiment decision: no split rule beat V6, V5 and V4 on both full means; retain existing submissions.


\subsection{Organizer feedback for V6 and local notebook audit}
The user's subsequent Kaggle screenshot reports V6 at $0.03833$, against
V5 at $0.01996$ and V4 at $0.01984$. V4 remains the best scored submission
among these versions; V6 is approximately $93.2\%$ worse than V4.
These organizer results do not support the local geometry gains transferring
to the organizer evaluation. All earlier local measurements remain recorded.

V6's displayed score equals V3's to five decimal places; this coincidence
does not establish a cause or prove a wrong upload. The local V6 notebook
matches its recorded artifact hashes, embeds \texttt{geo6\_threequarter}, and
its actual inference wrapper reproduces 500 synthetic action pairs against
the registered controller. This audit checks the local notebook only; the
uploaded Kaggle Version 6, its evaluation metadata and execution logs have
not been verified. V7's fixed local study remains exploratory, with no
organizer-superiority claim supported by a local pass alone.

\subsection{V7 rejection and return to V4: reheat-valve characteristics}
The V7 study completed 91 full episodes without rollout errors. Its only
dev-passing candidate, \texttt{geo7\_strong\_i}, reached dev mean
$0.0160358888$ ($-0.0853\%$ versus V6), but training mean $0.0142295906$
($+1.0743\%$ versus V6 and $+0.5119\%$ versus V5). It failed the predeclared
training gate; no V7 notebook was packaged. All three candidate outcomes and
earlier geometry measurements remain retained. V4 at $0.01984$ remains the
best known organizer score.

The next prospective study returns to the V4 controller without geometry
scaling. Its zone PI output $u$ commands damper $|u|$ and valve
$h=\max(0,-u)$. The simulator's coil heat depends on water and air capacity
rates and discharge-temperature limits; it is not assumed proportional to
the damper--valve product. The hypothesis is that valve response can be
adjusted independently of air delivery. Test $h^p$ with fixed
$p\in\{0.5,0.75,1.5\}$ during active PI, against the identity $p=1$ control.
Every curve preserves zero/full valve authority. The signed PI output,
integration condition, gains, damper, plant coordination and all V4 schedules
remain unchanged. Altered valve response may change dynamics and is not
assumed beneficial without complete rollouts.

After a fresh full identity control reproduces V4 within $10^{-10}$,
all three candidates receive all 20 dev scenarios. Dev-improving candidates
receive all 30 training scenarios in rank order until both full means improve
on V4. Only then are a notebook and standalone verification produced; the
bounded maximum is 151 new episodes before verification. Individual regressions
and component costs remain reported. Repeated local tuning is exploratory,
the test split remains unused, and only organizer evaluation can establish
that a submission beats V4.

New valve-curve notebooks print the variant name and a SHA-256 identifier of
the assembled UTF-8/LF policy source at startup, enabling an execution-log
comparison with the intended artifact. Full notebook-wrapper tests check
actions and this marker with a stub inference server; this is not a test of
the actual Kaggle service. Existing V4/V5/V6 artifacts remain unchanged.
Current tracked simulator files have no working-tree changes, but equality
with the organizer's installed SDK has not been established. Hash-frozen
inputs, memory/disk launch checks, checkpoints and atomic PDF replacement
accompany the new study.

\paragraph{Automatically recorded closed-loop results.}
Status: \texttt{complete}; stage: \texttt{finished}.
All entries below require complete 8,640-step episodes on matching scenarios.
\begin{center}\small
\begin{tabular}{llrrrr}\toprule
Cohort & Variant & $n$ & e\_sync & Candidate & Change \\\midrule
full dev & \texttt{valve8\_sqrt} & 20 & 0.01668440 & 0.01645237 & -1.39\% \\
full dev & \texttt{valve8\_gentle} & 20 & 0.01668440 & 0.01645037 & -1.40\% \\
full dev & \texttt{valve8\_soft} & 20 & 0.01668440 & 0.01740655 & +4.33\% \\
full training & \texttt{valve8\_gentle} & 30 & 0.01445394 & 0.01436451 & -0.62\% \\
\bottomrule\end{tabular}\end{center}
These are complete matched cohorts. Dev is repeatedly used model-selection data, not untouched validation. Individual regressions are reported without a hard cutoff.
Completed new rollouts: 91. No candidate replaces Version 4 automatically.
The unchanged-PI full-episode control score is $0.01620892$; exact-reproduction gate: passed.
The selected candidate is \texttt{valve8\_gentle}.


\subsection{Global-curve completion and conditional reheat follow-up}
The global valve-curve study completed 91 valid full episodes. Selected
\texttt{valve8\_gentle} ($p=0.75$) reached training mean $0.0143645138$
($-0.619\%$ versus V4) and dev mean $0.0164503660$ ($-1.403\%$).
It improves 14 training cases, regresses 15 and ties one; dev improves 12
and regresses eight. Energy/carbon reductions outweigh small increases in
mean comfort penalty. Standalone actions matched 1,152 real-observation pairs.
The user submitted this candidate as Kaggle notebook Version 7; its organizer
score is pending. Internal experiment labels do not denote Kaggle notebook
version numbers. The earlier local \texttt{geo7} rejection is a different study.
V4 remains the best known organizer-scored controller.

Weather-group summaries motivate but do not identify a mechanism: training
Fairbanks improves $2.742\%$ while training Phoenix worsens $1.122\%$, yet
dev Phoenix improves $3.004\%$. The next fixed study therefore uses observable
operating conditions, never a weather/building label. For unchanged V4 zone
demand $h=\max(0,-u)$, command valve opening
$h+w(h^{0.75}-h)$. Test (i) $w=1$ only when the final same-step AHU command is
running with positive normalized SAT; (ii)
$w=\operatorname{clip}((T_{\mathrm{low}}-T_{\mathrm{obs}})/0.5,0,1)$; and
(iii) the second weight only when the first command condition holds.
This command sign is a mode proxy, not measured supply-air temperature or
proof that the AHU adds heat. The inherited observation delay remains;
the policy does not read simulator state or reward.

The gate runs after parent coordination and hold logic, using that step's
final commands. Fixed-observation PI states/demands, damper and all plant
controls remain V4, with zero/full valve authority retained. Future feedback
may still change those commands through changed temperatures. Outside active
PI, V4 is inherited unchanged; no geometry dependence is introduced.

Two fresh full controls on \texttt{building\_train\_1/fairbanks\_1} must
reproduce V4 ($w=0$) and the submitted curve ($w=1$). All three gates then
receive all 20 dev scenarios. Candidates beating both full-dev reference
means receive all 30 training scenarios in rank order; packaging requires
both complete means below both references by more than $10^{-10}$, followed
by standalone verification. The maximum is 152 new episodes before those
checks. Regression details remain visible, repeated tuning is exploratory,
the test split stays unused, and organizer superiority requires a Kaggle score.
All preceding data and artifacts remain preserved.

\paragraph{Automatically recorded closed-loop results.}
Status: \texttt{complete}; stage: \texttt{finished}.
All entries below require complete 8,640-step episodes on matching scenarios.
\begin{center}\small
\begin{tabular}{llrrrr}\toprule
Cohort & Variant & $n$ & Curve & Candidate & Change \\\midrule
full dev & \texttt{valve9\_heat} & 20 & 0.01645037 & 0.01650237 & +0.32\% \\
full dev & \texttt{valve9\_cold} & 20 & 0.01645037 & 0.01667351 & +1.36\% \\
full dev & \texttt{valve9\_heat\_cold} & 20 & 0.01645037 & 0.01664084 & +1.16\% \\
\bottomrule\end{tabular}\end{center}
These are complete matched cohorts. Dev is repeatedly used model-selection data, not untouched validation. Individual regressions are reported without a hard cutoff.
Completed new rollouts: 62. No candidate replaces Version 4 automatically.
The unchanged-PI full-episode control score is $0.05873991$; exact-reproduction gate: passed.
The fresh global-curve reproduction gate passed.
Experiment decision: no gate beat both references on both full means; retain existing submissions.


\subsection{Organizer confirmation and execution supervision, 1 November}
The user's Kaggle screenshot confirms notebook Version 7
(\texttt{valve8\_gentle}) at $0.01922$, below V4's $0.01984$ by $3.125\%$.
It is now the best reported organizer score among these submissions. Earlier
pending/best-version statements record their historical state; all earlier
local and organizer results remain retained.

The conditional-curve experiment additionally receives independent background
supervision without changing its frozen policy or evaluation protocol. It
adopts a live scheduler, waits for surviving workers, validates frozen inputs
and resources before restarting from saved results, and logs scheduler exits
with bounded retry backoff. A real child-process crash/retry test complements
duplicate-launch and input/resource-gate tests. This improves execution
continuity, not controller performance. It does not prevent system sleep or
automatically start after reboot; those remain operational limitations.

\subsection{Conditional-curve rejection and Version 7 persistence study}
The conditional-curve study completed 62 valid episodes: two full reproduction
controls and 60 dev episodes. Relative to Version 7's global curve, the
command-mode, cold-zone and combined gates increased full-dev means by
$0.316\%$, $1.356\%$ and $1.158\%$. All failed the predeclared dev gate;
no training or notebook packaging followed. The preceding table and all
earlier results remain retained. Version 7 remains best at $0.01922$ publicly.

The next fixed study keeps that global $h^{0.75}$ valve curve and changes only
cooling-run persistence: 144, 432 or 576 five-minute steps (12, 36, 48 hours),
against the unchanged 288-step (24-hour) reference. The existing timer refreshes
on each new cooling request and cancels on a heating request; the duration
is persistence after the latest cooling call, not a total-run-length cap.
This differs from earlier off-window-only hold tests. Shorter persistence
may save ventilation, while longer persistence may preserve delivery;
both directions require measurement. Other PI, damper, plant and schedule
settings remain Version 7. No geometry or scenario identifiers choose actions.

A fresh full Version 7 control must reproduce its cached reward. All three
candidates receive all 20 dev episodes; eligible candidates then receive all
30 training episodes in rank order. Both complete means must beat Version 7
and V4 by more than $10^{-10}$ before packaging and standalone verification.
The maximum is 151 new full episodes before verification. Individual regressions
remain visible. Repeated local tuning is exploratory, the test split remains
unused, and only Kaggle establishes a score below $0.01922$. Independent
supervision and frozen-input checks accompany execution; older data remain intact.

\paragraph{Automatically recorded closed-loop results.}
Status: \texttt{prepared}; stage: \texttt{reference\_control}.
All entries below require complete 8,640-step episodes on matching scenarios.
\begin{center}\small
\begin{tabular}{llrrrr}\toprule
Cohort & Variant & $n$ & Curve & Candidate & Change \\\midrule
\bottomrule\end{tabular}\end{center}
These are complete matched cohorts. Dev is repeatedly used model-selection data, not untouched validation. Individual regressions are reported without a hard cutoff.
Completed new rollouts: 0. No candidate replaces Version 4 automatically.

